\section{PCA 从最大方差到最小重构误差}

设一家工厂在同一台设备上安装了几十个传感器。数据表虽然有几十列，许多列却会随着``负载''\,``温度''和``磨损''共同变化。直接把全部变量送进模型，不仅难以可视化，还可能把测量噪声当成新的方向。我们希望用少数坐标保留数据中最主要的结构。

但``主要''必须被定义。PCA 的答案是：选择能承载最大样本方差的正交方向；等价地，选择让正交投影重构误差最小的低维子空间。

\subsection{中心化不是可选的装饰}

把 \(n\) 个样本按行组成 \(X\in\mathbb R^{n\times p}\)，每行一个样本、每列一个变量，并令

\[
X_c=X-\mathbf 1\bar x^\top.
\]

其中 \(\mathbf 1\) 是全 1 列向量，\(\bar x\) 由各列的样本均值组成。若不中心化，原点到数据均值的距离也会被当作``变化''，第一方向可能只是指向均值，而不是描述样本之间怎样变化。PCA 分析的是围绕均值的结构，所以中心化决定了问题本身。

尺度选择则没有统一答案。若身高以米、收入以元记录，协方差 PCA 很可能被收入的数值尺度支配。标准化到单位方差相当于对相关矩阵做 PCA，给予各变量相近权重；但若变量单位相同且方差大小本就具有科学意义，机械标准化又会放大高噪声变量。尺度是建模假设，不是清洗步骤。

\subsection{最大方差怎样变成特征向量}

沿单位方向 \(v\) 的投影得分是 \(X_cv\)，其未归一化方差为

\[
\|X_cv\|_2^2
=v^\top X_c^\top X_cv.
\]

第一主方向因此解

\[
\max_{\|v\|_2=1}v^\top X_c^\top X_cv.
\]

这是 Rayleigh 商问题，最优 \(v_1\) 是 \(X_c^\top X_c\) 最大特征值对应的特征向量。后续方向在与前面方向正交的条件下依次最大化剩余方差。

若直接作 SVD：

\[
X_c=U\Sigma V^\top,
\]

则 \(V\) 的列是主方向，也称 loadings；\(X_cV=U\Sigma\) 是样本在新坐标下的 scores；第 \(j\) 个方向解释的样本方差为 \(\sigma_j^2/(n-1)\)。

一个最简单的例子是两个完全同步的中心化变量，协方差结构与

\[
\begin{pmatrix}
1&1\\
1&1
\end{pmatrix}
\]

成比例。方向 \((1,1)^\top/\sqrt2\) 承载全部变化，方向 \((1,-1)^\top/\sqrt2\) 上没有变化。原来的二维坐标其实只描述一个自由度。若第二个变量带少量独立噪声，第二奇异值不再为零，但仍远小于第一项。

\subsection{为什么最大方差等价于最小重构误差}

保留前 \(k\) 个方向组成 \(V_k\)，样本的低维表示为

\[
Z=X_cV_k,
\]

从该表示重构回原空间得到

\[
\widehat X_c=ZV_k^\top=X_cV_kV_k^\top.
\]

\(V_kV_k^\top\) 是到主子空间的正交投影。因为总平方长度可以分解为投影部分与正交残差部分，

\[
\|X_c\|_F^2
=\|X_cV_kV_k^\top\|_F^2
+\|X_c-X_cV_kV_k^\top\|_F^2,
\]

在总量固定时，最大化保留方差就等价于最小化平方重构误差。保留前 \(k\) 个奇异方向时，误差精确等于

\[
\sum_{j>k}\sigma_j^2.
\]

这不是 PCA 的偶然巧合：下一节的 Eckart--Young--Mirsky 定理将说明，没有任何秩不超过 \(k\) 的矩阵能在 Frobenius 范数下做得更好。

这也解释了 PCA 的边界：它最优的是线性子空间下的平方重构，不是``对任何下游任务都最优''。一个方差很小的方向可能恰好决定疾病类别；一个方差很大的方向也可能只是光照或批次效应。

\subsection{从公式走向建模流程}

如果 PCA 参与预测 pipeline，均值、尺度和主方向都只能由训练折估计，再原样应用于验证或测试数据。在全数据上先做 PCA 再交叉验证，会让验证样本影响表示，产生信息泄漏。

成分数 \(k\) 也不能只靠``累计解释方差达到 95\%''机械决定。应同时考察：

\begin{itemize}
\tightlist
\item
  scree plot 是否出现明显谱隙；
\item
  重构误差是否符合压缩或去噪需求；
\item
  下游验证性能是否在某个 \(k\) 后停止改善；
\item
  主子空间在不同抽样中是否稳定；
\item
  保留下来的方向能否结合变量语义解释。
\end{itemize}

单个 loading 的符号没有固有意义，因为 \(v\) 与 \(-v\) 表示同一轴；重复或非常接近的奇异值还会让该子空间内的具体基随扰动旋转，下一节讨论谱隙时会回到这种不稳定性。此时应解释整个子空间，而不是执着于某一列向量的方向。

离群点会因平方距离而强烈牵引主方向，非线性流形也可能无法被任何低维平面展开。PCA 没有失败到``算错''，而是优化了一个不再符合任务的目标。

\subsubsection{半开放任务}

构造两类二维数据：两类的共同变化方向方差很大，而区分类别的方向方差很小。比较只保留第一主成分与使用原始两维的分类性能。这个例子应能说明：为什么无监督重构最优不推出监督预测最优？然后把中心化操作错误地放到全数据上，观察交叉验证结果会怎样变化。

线性代数层面的完整推导见线性代数部分的``从数据矩阵到 PCA''一节。
